\subsection{From matrix entries to a kernel}\label{finite-kernel-model}

A matrix entry records how an operator moves one coordinate into another. A kernel plays a similar role when vectors are functions: its two variables describe an output location and an input location. We begin with a model that still contains only finitely many coordinates. Every identity will follow from finite-dimensional linear algebra and integrals whose convergence we can check directly.

Let \((X,\mu)\) be a sigma-finite measure space, and let \(H=L^2(X,\mu;\mathbb C)\). Thus vectors in \(H\) are square-integrable functions, with functions equal almost everywhere identified. Our inner product is
\[
\langle f,g\rangle=\int_X f(x)\overline{g(x)}\,d\mu(x).
\]
It is linear in the first argument and conjugate-linear in the second. Fix an orthonormal family \(u_1,\ldots,u_N\), so \(\langle u_i,u_j\rangle=\delta_{ij}\), and write \(V\) for its span. Choose measurable representatives of these functions. An empty family, with \(N=0\), is allowed.

Readers meeting this language for the first time may take \(X\) to be a finite set with counting measure. Functions are then column vectors, integrals are finite sums, and orthonormality is the familiar condition on vectors in \(\mathbb C^{|X|}\). The general setting also allows a finite family of functions on an interval or a surface.

For a complex matrix \(A=(a_{ij})\in M_N(\mathbb C)\), set
\[
K_A(x,y)=\sum_{i,j=1}^N a_{ij}u_i(x)\overline{u_j(y)},
\qquad
T_Af=\sum_{i,j=1}^N a_{ij}\langle f,u_j\rangle u_i.
\]
The first variable of \(K_A\) is the output variable. Indeed,
\[
(T_Af)(x)=\int_X K_A(x,y)f(y)\,d\mu(y)
\]
for almost every \(x\). To justify the integral, Cauchy--Schwarz gives
\[
\int_X |u_j(y)f(y)|\,d\mu(y)\leq\|u_j\|_2\|f\|_2=\|f\|_2.
\]
Only finitely many such terms occur. The displayed finite sum also defines a bounded operator on \(H\); for example,
\(\|T_Af\|_2\leq(\sum_{i,j}|a_{ij}|)\|f\|_2\).

The image of \(T_A\) lies in \(V\), and \(T_A\) vanishes on \(V^\perp\). Moreover,
\[
T_Au_j=\sum_i a_{ij}u_i.
\]
Consequently its restriction to \(V\) has matrix \(A\). Our integral operator is a finite matrix acting on \(V\), extended by zero on all remaining directions.

In particular, take \(X=\{1,\ldots,N\}\) with counting measure and \(u_j\) the indicator of the point \(j\). Then \(K_A(i,j)=a_{ij}\), and the integral formula is exactly ordinary matrix multiplication.

\subsection{A complete finite-rank trace identity}\label{finite-kernel-trace-identity}

Define the trace of this operator by
\(\operatorname{tr}(T_A)=\operatorname{tr}(T_A|_V)\).
The preceding section proves that the latter trace is independent of the basis of \(V\). Enlarging \(V\) by finitely many directions in \(V^\perp\) adds a zero block and changes nothing. Thus the unused zero directions do not affect this definition.

If the same operator has another such description on a finite-dimensional space \(W\), compare both descriptions inside \(V+W\). This space is invariant under the operator, and relative to \(V\oplus((V+W)\cap V^\perp)\) its matrix is the matrix on \(V\) together with a zero block. The same argument using \(W\) shows that the two traces agree.

\textbf{Proposition.} Let \(\lambda_1,\ldots,\lambda_N\) be the eigenvalues of \(A\), counted with algebraic multiplicity. Then the function \(x\mapsto K_A(x,x)\) given by the displayed kernel formula is integrable, and
\[
\boxed{\int_X K_A(x,x)\,d\mu(x)
=\operatorname{tr}(T_A)
=\sum_{j=1}^N\lambda_j.}
\]

\textbf{Proof.} Each diagonal term is integrable because
\[
\int_X |u_i(x)\overline{u_j(x)}|\,d\mu(x)
\leq\|u_i\|_2\|u_j\|_2=1.
\]
We may therefore integrate the finite sum term by term. Orthonormality gives
\[
\int_X K_A(x,x)\,d\mu(x)
=\sum_{i,j}a_{ij}\langle u_i,u_j\rangle
=\sum_i a_{ii}
=\operatorname{tr}(A).
\]
The matrix of \(T_A|_V\) is \(A\), and the trace--eigenvalue theorem of the preceding section gives the remaining equality. This includes \(N=0\), when all expressions vanish.

The eigenvalues here belong to the finite-dimensional restriction. We are not summing a list of all eigenvalues of an infinite-dimensional space, whose zero eigenspace could be infinite-dimensional. We have also made no assumption that \(A\) is diagonalizable.

In the useful special case \(A=\operatorname{diag}(\lambda_1,\ldots,\lambda_N)\), the kernel takes the simpler form
\[
K_A(x,y)=\sum_{j=1}^N\lambda_j u_j(x)\overline{u_j(y)}.
\]
Then \(T_Au_j=\lambda_j u_j\), and integrating its diagonal removes the normalized factors \(\int_X|u_j|^2\,d\mu=1\). This is a small but exact model of a spectral calculation.

\subsection{An explicit rank-one calculation}\label{finite-kernel-rank-one-example}

Take \(X=[0,1]\) with Lebesgue measure and \(u(x)=\sqrt{3}\,x\). Since \(\int_0^1 3x^2\,dx=1\), this is a unit vector. For \(\alpha\in\mathbb C\), the one-by-one matrix \((\alpha)\) gives
\[
K(x,y)=3\alpha xy,
\qquad
(Tf)(x)=3\alpha x\int_0^1 yf(y)\,dy.
\]
We have \(Tu=\alpha u\), while \(Tf=0\) whenever \(f\perp u\). Directly,
\[
\int_0^1 K(x,x)\,dx
=3\alpha\int_0^1x^2\,dx=\alpha.
\]
The calculation remains valid for \(\alpha=0\); then the operator has rank zero rather than one.

\subsection{Why the specified kernel matters}\label{kernel-diagonal-representatives}

Our diagonal is defined by the finite formula for \(K_A\). Changing the representatives of the finitely many \(u_j\) changes that diagonal only on a null set, so its integral is well defined.

An arbitrary representative of a kernel defined only almost everywhere on \(X\times X\) need not have this property. On \([0,1]^2\), the zero function and the function that is one on \(x=y\) and zero elsewhere agree almost everywhere for two-dimensional Lebesgue measure. Both give the zero integral operator. Their diagonal integrals are nevertheless zero and one, respectively. The diagonal itself has product measure zero.

Thus an integral operator alone does not justify taking the diagonal of an arbitrary measurable kernel. Our finite separated formula supplies a valid diagonal. Later, continuity, smoothness, convergence estimates, or a specified regularization must provide the corresponding justification. The phrase ``integrate the diagonal'' always carries this mathematical obligation.

\subsection{Exercises with solutions}\label{finite-kernel-exercises}

\textbf{Exercise 1.} Set \(A=I_N\). Identify \(T_A\) and compute its trace without finding a characteristic polynomial.

\textbf{Solution.} The formula becomes \(T_Af=\sum_j\langle f,u_j\rangle u_j\), the orthogonal projection onto \(V\): its difference from \(f\) is orthogonal to every \(u_j\). It fixes each of the \(N\) basis vectors and kills \(V^\perp\), so its trace is \(N\). Equivalently, its diagonal integral is \(\sum_j\|u_j\|_2^2=N\).

\textbf{Exercise 2.} For any \(v\in H\), consider \(Sf=\langle f,v\rangle v\). What replaces the normalization factor one in the trace calculation?

\textbf{Solution.} If \(v=0\), then \(S=0\). Otherwise set \(u=v/\|v\|_2\). We have \(Sf=\|v\|_2^2\langle f,u\rangle u\), so the trace is \(\|v\|_2^2\). The kernel \(v(x)\overline{v(y)}\) has diagonal integral \(\int_X|v|^2\,d\mu\), giving the same answer. Normalization is part of the formula, not a cosmetic choice.
